%% ---------------------------------------------------------------------
\section{Fusing genuinely different sources}\label{sec:multisource}

Everything to this point fuses one knowledge source into one image stream.
This section asks whether the same rule governs fusion in the sense the
term usually carries: two \emph{sensors}, and two \emph{classifiers}.
Section~\ref{sec:multisensor} takes visible and non-visible Sentinel-2
bands as separate sources; Section~\ref{sec:decisionfusion} fuses trained
models at the decision level. Both are tests of transfer rather than
restatements: the rule was derived from training-time fusion of a prior
into a network, and neither of these is that.

\subsection{Two sensors: multispectral EuroSAT}\label{sec:multisensor}
EuroSAT ships as 13 Sentinel-2 bands, of which three are visible. The
visible bands and the red-edge, near-infrared, short-wave infrared and
cirrus bands come from different detectors and carry different physical
information, which makes the complementary-versus-redundant question
askable of the data itself. We build three sources from one archive:
\textbf{visible} (3 bands), \textbf{non-visible} (the other 10), and
\textbf{sensor-fused} (all 13). The tiles and the committed subset indices
are identical across the three, so the only thing that varies is which
detectors the network reads. Each is crossed with the baseline and the
prior under the frozen recipe.

\begin{table}[t]
\caption{Multispectral EuroSAT. Identical tiles and identical subset
indices; only the band set differs. Three seeds per cell, ResNet-18. The
prior's gain \emph{grows} with the number of sensors present rather than
being absorbed by them, and the same ordering appears in the frozen-feature
evaluation, so it is a feature-level effect.}
\label{tab:sensorfusion}
\centering\scriptsize
\setlength{\tabcolsep}{2.6pt}
\zebra{2}
\begin{tabular}{@{}P{0.26\columnwidth}rrrr@{}}
\toprule
Source & Baseline~\up & + Prior~\up & $\Delta$~\up & $G$~\up \\
\midrule
\multicolumn{5}{@{}l}{\emph{5\% of the training split (1{,}080 tiles)}} \\
Visible (3 bands) & 89.56 & 91.19 & $+1.63$ & -- \\
Non-visible (10) & 90.14 & 92.73 & $+2.60$ & -- \\
Sensor-fused (13) & \textbf{91.48} & \textbf{94.78} & $\mathbf{+3.30}$ & -- \\
\midrule
\multicolumn{5}{@{}l}{\emph{25\% of the training split}} \\
Visible (3 bands) & 97.88 & 97.85 & $-0.02$ & $-0.02$ \\
Non-visible (10) & 97.77 & 98.19 & $+0.43$ & $+0.28$ \\
Sensor-fused (13) & \textbf{98.33} & \textbf{98.65} & $+0.33$ & $+0.37$ \\
\bottomrule
\end{tabular}
\end{table}

Table~\ref{tab:sensorfusion} reports three things.

\emph{The two band sets are complementary.} Fusing them beats the better
single sensor by $+1.34$ at 5\%, $+1.15$ at 10\% and $+0.45$ at 25\%. This
is the stack outcome, measured on the data rather than on a training
intervention, and it decays with data exactly as every other effect in this
study does: by 25\% of a 21{,}600-tile archive there is little left for a
second sensor to contribute.

\emph{The prior's gain grows with the number of sensors.} At 5\% it is
$+1.63$ on the visible bands, $+2.60$ on the non-visible, and $+3.30$ on
both. Had oriented-energy structure been a substitute for spectral
coverage, adding detectors would have absorbed it; instead the two
compound. This is the sharpest available statement of the currency account,
because the second source here is a physical sensor rather than another
training trick, and its currency is unambiguously not the prior's.

\emph{And it is a feature-level effect.} Under identical frozen-feature
evaluation, $G$ moves in the same order as $\Delta$: $-0.02$, $+0.28$,
$+0.37$. The law's arithmetic also holds where it should. All three
baselines sit near 98\%, far above the crossing bracket, where the account
predicts a readout term of ${\approx}0$; the measured residuals are
$0.00$, $+0.15$ and $-0.04$.

\subsection{Two classifiers: decision-level fusion}\label{sec:decisionfusion}
Multi-classifier fusion is a different level again, and the rule was not
derived there. Using trained checkpoints from the CIFAR-100 5\% cells and
no further training, we average class posteriors for every pair of arms and
record the gain over the better single arm together with the pair's
disagreement rate, the classical diversity measure
\citep{brown2005diversity}.

The first thing the data say is a caution rather than a confirmation.
Ensemble gain is well explained ($R^2 = 0.74$) by two terms, and the
dominant one is not diversity:
\begin{equation}
  \mathrm{gain} \;=\; -2.74 \;+\; 0.089\,\mathrm{disagreement}
                       \;-\; 0.341\,|\Delta_{\mathrm{acc}}| ,
  \label{eq:ensgain}
\end{equation}
so pairing a strong arm with a weak one destroys the gain regardless of
what either contributes. Restricting to pairs within one point of each
other, where currency is actually the free variable, leaves three:

\begin{center}\scriptsize
\begin{tabular}{@{}P{0.34\columnwidth}rr@{}}
\toprule
Pair & Gain~\up & Disagreement \\
\midrule
Prior + SimCLR (50 ep) & $+1.20$ & 47.2\% \\
SimCLR (50 ep) + prior, wide bank & $+1.13$ & 47.2\% \\
Prior + prior, wide bank & $-0.04$ & 28.1\% \\
\bottomrule
\end{tabular}
\end{center}

Two models built on the same knowledge source disagree on 28\% of test
images and ensembling them buys nothing; pair either with a
self-supervised model and disagreement rises to 47\% and ensembling buys
${\sim}1.2$ points. That is the currency rule with its mechanism visible,
and it is what the classical account would predict.

We report the qualification with equal weight, because it bounds the
taxonomy rather than extending it. The prior and SimCLR do \emph{not} stack
at training time (Section~\ref{sec:substitute}), yet they ensemble usefully
at decision time ($+1.24$). Redundancy in the sense of ``jointly training on
both adds nothing'' therefore does not imply redundancy in the sense of
``their errors are correlated''. The substitution result is a statement
about training-time fusion, and we do not claim it transfers to decision
fusion. With three matched pairs this is suggestive rather than
established, and we say so.

\subsection{A procedure for deciding whether to fuse}\label{sec:procedure}
The practical content of the taxonomy is that the outcome of a combination
can be estimated \emph{before} the combined system is built, from
measurements of each source alone. We state that as a procedure rather than
as a finding, since it is what a practitioner would actually run.

\begin{table}[t]
\caption{Deciding whether to fuse two sources, before training the
combination. Every quantity is measured on the sources separately; the only
cost beyond training each source once is two frozen-feature evaluations.}
\label{tab:procedure}
\centering\scriptsize
\setlength{\tabcolsep}{2.6pt}
\zebra{2}
\begin{tabular}{@{}P{0.05\columnwidth}P{0.86\columnwidth}@{}}
\toprule
\multicolumn{2}{@{}l}{\textbf{Inputs}: sources $A$, $B$; a baseline arm; a labelled split} \\
\midrule
1 & Train the baseline, $A$ alone and $B$ alone. Record each arm's
    end-to-end gain and declared compute multiple. \\
2 & Evaluate all three frozen representations with an identical linear
    protocol. This gives $G_A$ and $G_B$, each source's feature-level
    contribution. \\
3 & If either source's gain is ${\approx}0$, it supplies nothing to
    overlap: the other source's gain is available in full, and the
    combination is not the question. \\
4 & If $G_A \approx G_B$ \emph{and} both arms improve the same cells,
    the sources carry the same currency. Expect \textbf{substitution}:
    the combination will not exceed the better single arm. Take the
    cheaper source. \\
5 & If $G_A$ and $G_B$ differ in kind, expect \textbf{stacking}, and a
    combined gain larger than either alone. \\
6 & If one source is a mature initialization, do not inject a shaping
    prior on top of it. Expect a \textbf{tax} proportional to what the
    initialization was contributing, measurable in advance as that
    source's advantage over training from scratch. \\
\bottomrule
\end{tabular}
\end{table}

Table~\ref{tab:procedure} states it. Three properties are worth drawing
out. It is \emph{cheap}: the only cost beyond training each source once is
two linear evaluations on frozen features, which run in minutes on cells
that took hours to train. It is \emph{predictive rather than
retrospective}, which places it alongside transferability estimation
\citep{nguyen2020leep,you2021logme}, with the difference that those methods
score a single pre-trained source for a target task while this scores a
\emph{pair} for combination. And it is \emph{falsifiable}: step 4 asserts
that equal feature gains imply no benefit from combining, which
Sections~\ref{sec:substitute} and \ref{sec:decisionfusion} both test and
which the training-time-versus-decision-time discrepancy above already
qualifies.

The honest limit is step 4's antecedent. We compare feature gains by
magnitude and by which cells improve, not by a formal decomposition of the
information each source carries. Partial information decomposition
\citep{williams2010nonnegative,liang2023quantifying} is the principled
version of that comparison, and an estimator that scaled to these
representations would replace our proxy with a quantity that means
something on its own.
